\section{Conclusion}
\label{sec:conclusion}

%==============================================================================
% SKELETON ONLY.  Status: WAIT_SIM_REAL.
% The conclusion of a predictivity study is determined by the predictivity
% data.  Writing it now would fix an answer the experiments have not given.
%==============================================================================
% [1] Restate what was measured (not what was built): predictivity of an
%     underwater simulator for camera-driven obstacle avoidance, under a
%     progressive measured-calibration ladder.
% [2] Answer RQ1: how close simulated safety/trajectory metrics were to
%     physical ones .................................................. \SREAL
% [3] Answer RQ2: which calibration class changed predictivity, and which
%     did not (a null at some level is a result) ..................... \SREAL
% [4] Answer RQ3: whether the ranking of avoidance stacks survived
%     deployment ..................................................... \SREAL
% [5] Answer RQ4 (already available): temporal persistence predicts better
%     through maneuver onset; the Kalman filter's gating provides outlier
%     robustness --- two separate benefits, not one.
% [6] Consequence for practice: what a marine roboticist should calibrate
%     first if they intend to use a simulator to CHOOSE between avoidance
%     methods.  This is the sentence the paper exists to deliver, and it
%     cannot be written before [3].
% [7] Next step: open release of code, configurations and experiment
%     manifests.
\PENDING{CONCLUSION --- structure fixed above; determined by the sim-to-real
data, not before.}
